\begin{textsample}{POS Dim 2 – 2011\_c4 – Score 10.00 – 2011\_c4/enade\_2011\_eng\_g\_5\_6b7bd07019.txt}  \label{ex:f2_pos_122}
CAPÍTULO 4 PERCEPÇÃO DA PROVA As análises feitas neste capítulo tratam das percepções dos concluintes da Área de Engenharia – Grupo V sobre a prova aplicada no ENADE/2011.
Estas percepções foram mensuradas por meio de nove questões que avaliaram desde o grau de dificuldade da prova até o tempo gasto para concluí -la.
As percepções sobre a prova foram relacionadas com o desempenho dos estudantes e com a Grande Região de funcionamento do curso.
O \textbf{questionário} de percepção da prova encontra -se ao final do Anexo IV, que traz a \textbf{reprodução} da prova.
O desempenho dos estudantes foi classificado em quatro quartos.
Para tanto, esse desempenho foi \textbf{ordenado} de forma ascendente.
O percentil 25, P25, também conhecido como primeiro quartil, é a nota de desempenho que deixa um quarto ( 25 \% ) dos valores observados abaixo e três quartos acima.
A Figura 1 apresenta uma ilustração deste conceito.
O quarto inferior de desempenho é composto pelas notas abaixo do primeiro quartil.
Já o percentil 75, P75, também conhecido como terceiro quartil, é o valor para o qual há três quartos ( 75 \% ) dos dados abaixo e um quarto acima dele.
O quarto superior de desempenho é composto pelas notas iguais ou acima do terceiro quartil.
O percentil 50, P50, também conhecido como mediana, é o valor que divide as notas em dois conjuntos de igual tamanho.
O segundo quarto inclui valores entre o primeiro quartil ( P25 ) e a mediana.
O terceiro quarto contém os valores entre a mediana ( P50 ) e o terceiro quartil ( P75 ).
Vale ressaltar que percentis, quartis e medianas são pontos que não obrigatoriamente pertencem ao conjunto original de dados, ao passo que os quartos são \textbf{subconjuntos} dos dados originais.
A seguir, serão apresentados gráficos com resultados selecionados, relativos às nove questões avaliadas por grupos de estudantes.
Os gráficos apresentam nas barras o percentual de alunos que assinalaram uma das opções ou a soma das porcentagens daqueles que assinalaram duas ( ou três ) delas.
Por exemplo, para as questões 1 e 2, os gráficos apresentam a porcentagem total de participantes que assinalaram as opções ( D ) difícil e ( E ) muito difícil.
Em cada barra foram assinalados também os extremos do intervalo de confiança de 95 \% como linhas verticais unidas por uma linha horizontal na forma da letra H maiúscula, semelhantemente aos gráficos do Capítulo 3.
As Tabelas no Anexo II apresentam os valores absolutos e a distribuição percentual das alternativas válidas das nove questões, segundo o mesmo \textbf{recorte} de desempenho dos alunos e Grande Região de funcionamento do curso. 4.1 GRAU DE DIFICULDADE DA PROVA 4.1.1 Componente de Formação Geral Ao avaliarem “ Qual o grau de dificuldade desta prova na parte de Formação Geral? ” ( Questão 1 ), 16,2 \% do grupo de inscritos e presentes optaram pelas alternativas difícil ou muito difícil.
Entretanto, para mais da metade dos estudantes ( 59,6 \% ), o Componente de Formação Geral da prova foi considerado com grau de dificuldade médio ( Gráfico 4.1, Gráfico 4.2 e, no Anexo II, a Tabela II.1 ).
A região Centro-Oeste não apresentou inscritos.
O percentual de estudantes que consideraram a prova como difícil ou muito difícil foi maior na região Sul, onde a proporção foi de 19,9 \%, enquanto a de menor incidência foi a Norte, com 14,3 \%.
No Gráfico 4.1 é possível observar que esta diferença não é estatisticamente significativa.
Nas Grandes Regiões, a proporção de presentes à prova que consideraram o Componente de Formação Geral como sendo de grau de dificuldade médio esteve entre 57,5 \% na região Sudeste e 64,0 \% na região Sul.
O percentual de alunos que consideraram a prova difícil ou muito difícil apresentou tendência decrescente conforme o aumento do desempenho, variando de 28,8 \% ( 1º quarto ) até 6,4 \% ( 4º quarto ).
A alternativa modal para esta pergunta foi médio, com respectivamente 60,1 \% e 59,1 \% dos respondentes destes quartos extremos. 4.1.2 Componente de Conhecimento Específico Ao responderem à Questão 2 – “ Qual o grau de dificuldade desta prova na parte de Componente Específico? ” – 36,2 \% do grupo de estudantes classificaram -na como difícil ou muito difícil.
Além disso, o Componente de Conhecimento Específico da prova foi considerado com grau de dificuldade médio por 58,9 \% dos alunos ( Gráfico 4.3, Gráfico 4.4, e, no Anexo II, a Tabela II.2 ).
A análise das respostas dos estudantes quanto ao grau de dificuldade do Componente de Conhecimento Específico da prova, agregado por Grande Região, mostra que a diferença entre a maior e a menor proporção de alunos que a avaliaram como difícil ou muito difícil não é estatisticamente significativa : 45,7 \% na região Norte e 34,4 \% na Sudeste.
O percentual de alunos que classificaram o grau de dificuldade como médio, no Componente de Conhecimento Específico, variou de 51,4 \% a 60,4 \%, para as mesmas regiões, Norte e Sudeste, respectivamente.
Considerando -se a avaliação da dificuldade das questões do Componente de Conhecimento Específico da prova, de acordo com o desempenho dos estudantes, não se observa diferença estatisticamente significativa, a exceção do 1º com o 4º quarto, mas uma tendência decrescente é discernível.
Em todos os quartos, a proporção dos que classificaram a parte específica como difícil ou muito difícil variou de 26,5 \% ( 4º quarto ) a 43,7 \% ( 1º quarto ).
A alternativa modal para a Questão 2 foi o médio, com 52,3 \% do quarto inferior e 68,1 \% do superior optando por esta resposta. 4.2 EXTENSÃO DA PROVA EM RELAÇÃO AO TEMPO TOTAL \textbf{Indagados} quanto à extensão da prova, em relação ao tempo total oferecido para a sua resolução ( Questão 3 ), os estudantes apontaram, com maior incidência, a alternativa que considerava a extensão adequada, para todas as agregações consideradas, a exceção da região Nordeste ( Gráfico 4.5, Gráfico 4.6, e, no Anexo II, a Tabela II.3 ), cujos estudantes optaram por longa numa proporção de 44,9 \%.
O percentual de alunos que responderam ser a extensão da prova adequada foi de 60,1 \%.
Já 35,2 \% dos inscritos presentes consideraram que o exame foi longo ou muito longo e menos do que 5 \% o avaliaram como curto ou muito curto.
Entre as Grandes Regiões a proporção daqueles que avaliaram a prova como longa ou muito longa em relação ao tempo total destinado à sua resolução variou bastante : de 22,9 \% na região Norte até 63,8 \% na região Nordeste.
As diferenças entre a região Norte e as demais são estatisticamente significativas.
Considerando -se o desempenho dos alunos, dentre os que optaram pela extensão da prova como adequada, apresentou uma tendência crescente a partir do 2º quarto, variando de 51,9 \% até 65,5 \%.
No Gráfico 4.6, pode -se constatar que para os quatro quartos de desempenho a proporção de estudantes que consideraram a prova longa ou muito longa, também apresenta uma tendência decrescente conforme o aumento do desempenho a partir do 2º quarto, sem diferenças estatisticamente significativas. 4.3 COMPREENSÃO DOS ENUNCIADOS DAS QUESTÕES 4.3.1 Componente de Formação Geral Com relação aos enunciados das questões do Componente de Formação Geral ( Questão 4 ), as opiniões foram positivas, já que 82,2 \% dos alunos avaliados consideraram os enunciados de todas ou da maioria das questões claros e objetivos ( Gráfico 4.7, Gráfico 4.8, e, no Anexo II, a Tabela II.4 ).
Na análise regional, a \textbf{percentagem} de estudantes que avaliaram que todos ou a maioria dos enunciados das questões do Componente de Formação Geral estavam claros e objetivos variou de 75,4 \% na região Sul a 84,4 \% na região Sudeste, não havendo diferenças estatisticamente significativas entre as Grandes Regiões.
A análise das percepções dos estudantes sobre a clareza e objetividade dos enunciados permite afirmar que todos ou a maioria dos enunciados de questões relativas ao Componente de Formação Geral foram considerados claros e objetivos para a maior parte dos respondentes ( maior do que 75 \% em todas as regiões e para todos os quartos de desempenho ).
Segundo o desempenho, observa -se que a proporção dos que emitiram esta opinião \textbf{cresce}, \textbf{grosso} modo, do 1º quarto de desempenho para o 4º quarto, sem diferenças estatisticamente significativas.
No quarto superior, a clareza e objetividade de todos ou da maioria dos enunciados das questões foi percebida por 85,3 \% dos alunos e no quarto de desempenho inferior tal avaliação foi emitida por 76,3 \% deles. 4.3.2 Componente de Conhecimento Específico Com relação aos enunciados das questões do Componente de Conhecimento Específico da prova, para 82,2 \% dos estudantes avaliados da Área de Engenharia – Grupo V a clareza e a objetividade ( Questão 5 ) estavam presentes em todas ou na maioria das questões ( Gráfico 4.9, Gráfico 4.10, e no Anexo II, a Tabela II.5 ).
A maioria dos estudantes de todas as Grandes Regiões brasileiras considerou claros e objetivos todos ou a maioria dos enunciados das questões do Componente de Conhecimento Específico da prova, percentual sempre maior do que 75 \%.
A diferença entre 75,9 \% ( Sul ) e 88,6 \% ( Norte ) não é estatisticamente significativa.
A proporção de estudantes que consideraram os enunciados das questões claros e objetivos apresenta uma tendência decrescente em relação ao aumento de desempenho a partir do 2º quarto : mais elevada no 2º quarto ( 87,3 \% ) se comparada ao quarto superior de desempenho ( 82,4 \% ).
Esta diferença não é estatisticamente significativa. 4.4 SUFICIÊNCIA DAS INFORMAÇÕES/INSTRUÇÕES FORNECIDAS Ao avaliarem as informações/instruções fornecidas para a resolução das questões ( Questão 6 ), 92,3 \% dos respondentes da Área de Engenharia – Grupo V de todo o Brasil afirmaram que estas eram até excessivas ou suficientes em todas ou na maioria das questões ( Gráfico 4.11, Gráfico 4.12, e, no Anexo II, a Tabela II.6 ).
Quanto à distribuição de respondentes pelas Grandes Regiões observa -se que a proporção de estudantes que consideraram as informações/instruções fornecidas até excessivas ou suficientes em todas ou na maioria das questões foi sempre superior a 85 \%, chegando a 94,2 \% na região Nordeste.
As diferenças entre as regiões não chegaram a ser estatisticamente significativas.
Levando -se em conta o desempenho dos participantes, notam -se diferenças estatisticamente significativas entre as opiniões de estudantes do quarto inferior e superior de desempenho, como mostra o Gráfico 4.12.
O percentual de participantes que avaliaram as informações/instruções como até excessivas ou suficientes em todas ou na maioria das questões foi mais elevado no quarto superior ( 96,1 \% ), percentual superior à média nacional ( 92,3 \% ).
Já no quarto inferior, a suficiência das informações/instruções declarada como até excessiva em todas ou na maioria das questões foi percebida por 87,3 \% dos respondentes. 4.5 DIFICULDADE ENCONTRADA AO RESPONDER À PROVA Perguntados sobre as dificuldades com as quais se depararam ao responder à prova ( Questão 7 ), 18,0 \% dos estudantes apontaram o desconhecimento do conteúdo.
Para 33,3 \%, a forma diferente de abordagem do conteúdo foi indicada como dificuldade.
Já a falta de motivação para fazer a prova foi a dificuldade apontada por 23,6 \% dos respondentes.
Considerando -se todo o Brasil, 20,1 \% dos respondentes afirmaram que não tiveram qualquer tipo de dificuldade para responder à prova ( Tabela II.7 no Anexo II ).
Os Gráficos 4.13 e 4.14 apresentam os percentuais de estudantes que apontaram o desconhecimento do conteúdo como dificuldade percebida ao responder à prova.
Na análise por Grandes Regiões, o percentual de inscritos e presentes que apontaram o desconhecimento do conteúdo como dificuldade ao responder à prova não superou 28 \%.
Os percentuais variaram de 14,3 \% na região Norte a 27,9 \% na Nordeste, não sendo esta diferença estatisticamente significativa.
A forma diferente de abordagem do conteúdo foi a escolha modal em três regiões.
Os percentuais desta escolha variaram de 28,6 \% ( região Norte – nesta região a escolha modal foi falta de motivação ) a 33,9 \% ( Sul ).
O percentual de alunos que citou a falta de motivação como dificuldade variou de 14,7 \% ( região Nordeste ) a 34,3 \% ( região Norte ).
Os que declararam não ter qualquer dificuldade para responder à prova variaram de 13,4 \% na região Sul a 23,1 \% na Sudeste.
Com relação aos quartos de desempenho, o desconhecimento do conteúdo, foi a opção escolhida por 16,9 \% dos estudantes do quarto superior e 18,9 \% do quarto inferior.
A alternativa modal para os alunos, quando agregados pelos quartos de desempenho, foi que a dificuldade encontrada foi causada pela forma diferente de abordagem do conteúdo : 37,8 \% no quarto inferior e 27,9 \% do quarto superior assim o responderam. 4.6 CONTEÚDOS DAS QUESTÕES OBJETIVAS DA PROVA Ao analisarem os conteúdos das questões objetivas da prova ( Questão 8 ), um percentual muito pequeno dos estudantes avaliados, apenas 4,6 \%, afirmou que não estudou ainda a maioria desses conteúdos ( Gráficos 4.15, Gráfico 4.16, e a Tabelas II.8 no Anexo II ).
A maioria ( 78,2 \% ) afirmou ter estudado e aprendido muitos ou todos os conteúdos.
Na análise por Grande Região, a proporção de respondentes que escolheram a opção não estudou ainda a maioria desses conteúdos, foi pequena.
No entanto, a região Sul apresenta um percentual de respondentes muito acima da média praticada pelas demais regiões ( 11,2 \% ).
A região Sul apresenta diferenças estatisticamente significativas na comparação com as regiões Nordeste e Sudeste.
Em todas as regiões, a maioria dos presentes afirmou ter estudado e aprendido muitos ou todos os conteúdos, com proporções variando entre 57,1 \% na região Norte e 83,6 \% na Sudeste.
Considerando -se separadamente as opiniões de estudantes dos quatro quartos de desempenho, observa -se que, no quarto inferior, 10,1 \% ofereceram como resposta que não estudou ainda a maioria desses conteúdos, sendo 2,5 \% os do quarto superior com a mesma resposta.
A diferença entre os alunos que optaram por este motivo de dificuldade nos quartos extremos é estatisticamente significativa.
Tendo em conta o quarto superior, 90,2 \% dos alunos afirmaram ter estudado e aprendido muitos ou todos os conteúdos. 4.7 TEMPO GASTO PARA CONCLUIR A PROVA Ao responderem sobre o tempo de conclusão da prova ( Questão 9 ), mais de quatro quintos dos estudantes ( 81,2 \% ) afirmou ter gasto entre duas e quatro horas ( Gráfico 4.17, Gráfico 4.18 e, no Anexo II, a Tabela II.9 ).
Considerando -se as Grandes Regiões brasileiras, os que utilizaram entre duas e quatro horas para finalizar a prova ficaram abaixo do percentual nacional na região Norte ( 71,4 \% ).
Nas demais regiões, Nordeste, Sudeste e Sul, o percentual de alunos que \textbf{dispensaram} entre duas e quatro horas para concluir a prova, ficou em torno de 82 \%, como mostra o Gráfico 4.17.
Uma vez sendo analisadas as alternativas escolhidas pelos estudantes que se situam nos diferentes quartos de desempenho, observa -se que uma maior proporção de participantes no quarto superior declarou ter gasto entre duas e quatro horas para concluir a prova quando comparados com os do quarto inferior, respectivamente 89,6 \% e 73,2 \%.
As diferenças entre os três primeiros quartos não são estatisticamente significativas e as proporções no gráfico evidenciam uma tendência crescente.

% matched lemmas: crescer, dispensar, grosso, indagar, ordenar, percentagem, questionário, recorte, reprodução, subconjunto
\end{textsample}
